% !TeX program = lualatex
% !TeX root = predavanje07.tex
% ============================================================
% Zajednički preamble za sva predavanja RWA kolegija
% Uključiti s: % ============================================================
% Zajednički preamble za sva predavanja RWA kolegija
% Uključiti s: % ============================================================
% Zajednički preamble za sva predavanja RWA kolegija
% Uključiti s: \input{../common/preamble.tex}
% ============================================================

\documentclass[aspectratio=149,11pt]{beamer}
\usetheme{CambridgeUS}

\usepackage{fontspec}              % radi samo sa XeLaTeX/LuaLaTeX
\setsansfont{Latin Modern Sans}    % sans-serif font
\setmonofont{Latin Modern Mono}    % monospace font za kod
\usepackage[croatian]{babel}
\usepackage{csquotes}              % \enquote{} → hrvatski navodnici „..."
\usepackage{lmodern}
\usepackage{booktabs}
\usepackage{listings}
\usepackage{xcolor}
\usepackage{hyperref}

% ---------- Autor i logo ----------
\author{mag.ing. Josip Torić (SIT UNIZD)}
\titlegraphic{\includegraphics[height=2.5cm]{../common/SIT_logo_tekst1.pdf}}

% ---------- Listings ----------
\definecolor{codebg}{RGB}{248,248,248}
\definecolor{codeframeclr}{RGB}{220,220,220}
\lstdefinelanguage{bash}{
  morekeywords={sudo,apt,brew,pip,pipx,uv,winget,python,git,flask,source,deactivate,code,gunicorn,docker,compose,heroku,render,railway,node,npm,pnpm,yarn,uvicorn,alembic,pytest,ruff,black,precommit},
  sensitive=false,
  morecomment=[l]{\#}
}
\lstdefinelanguage{python}{
  morekeywords={from,import,def,class,return,if,elif,else,for,in,while,try,except,raise,with,as,not,and,or,is,lambda,yield,True,False,None,await,async},
  sensitive=true,
  morecomment=[l]{\#},
  morestring=[b]',
  morestring=[b]"
}
\lstdefinelanguage{sql}{
  morekeywords={SELECT,FROM,WHERE,INSERT,INTO,VALUES,UPDATE,SET,DELETE,CREATE,TABLE,PRIMARY,KEY,FOREIGN,REFERENCES,UNIQUE,NOT,NULL,INDEX,ON,ALTER,DROP,CONSTRAINT},
  sensitive=false,
  morecomment=[l]{--}
}
\lstdefinelanguage{json}{
  morecomment=[l]{//},
  morestring=[b]",
  showstringspaces=false
}
\lstdefinelanguage{typescript}{
  morekeywords={import,export,from,const,let,var,function,return,if,else,for,of,in,while,switch,case,break,async,await,try,catch,finally,throw,new,class,this,extends,implements,interface,type,enum,readonly,public,private,protected,static,get,set,true,false,null,undefined,void,string,number,boolean,any,unknown,never,typeof,keyof,as},
  sensitive=true,
  morecomment=[l]{//},
  morecomment=[s]{/*}{*/},
  morestring=[b]',
  morestring=[b]",
  morestring=[b]`
}
\lstalias{javascript}{typescript}
\lstset{
  backgroundcolor=\color{codebg},
  basicstyle=\ttfamily\footnotesize,
  frame=single,
  rulecolor=\color{codeframeclr},
  frameround=tttt,
  breaklines=true,
  tabsize=2,
  showstringspaces=false,
  upquote=true,
  columns=fullflexible
}

% ---------- Helper macros ----------
\newcommand{\codefile}[1]{\texttt{#1}}
\newcommand{\kbd}[1]{\texttt{#1}}

% ---------- Code frame environment ----------
% Automatski dodaje fragile i shrink=5, koristi se umjesto \begin{frame}
% Primjer: \begin{codeframe}{Naslov} ... \end{codeframe}
% Primjer s opcijama: \begin{codeframe}[shrink=15]{Naslov} ... \end{codeframe}
\newenvironment{codeframe}[2][]
  {\begin{frame}[fragile,shrink=5,#1]{#2}}
  {\end{frame}}

% ============================================================

\documentclass[aspectratio=149,11pt]{beamer}
\usetheme{CambridgeUS}

\usepackage{fontspec}              % radi samo sa XeLaTeX/LuaLaTeX
\setsansfont{Latin Modern Sans}    % sans-serif font
\setmonofont{Latin Modern Mono}    % monospace font za kod
\usepackage[croatian]{babel}
\usepackage{csquotes}              % \enquote{} → hrvatski navodnici „..."
\usepackage{lmodern}
\usepackage{booktabs}
\usepackage{listings}
\usepackage{xcolor}
\usepackage{hyperref}

% ---------- Autor i logo ----------
\author{mag.ing. Josip Torić (SIT UNIZD)}
\titlegraphic{\includegraphics[height=2.5cm]{../common/SIT_logo_tekst1.pdf}}

% ---------- Listings ----------
\definecolor{codebg}{RGB}{248,248,248}
\definecolor{codeframeclr}{RGB}{220,220,220}
\lstdefinelanguage{bash}{
  morekeywords={sudo,apt,brew,pip,pipx,uv,winget,python,git,flask,source,deactivate,code,gunicorn,docker,compose,heroku,render,railway,node,npm,pnpm,yarn,uvicorn,alembic,pytest,ruff,black,precommit},
  sensitive=false,
  morecomment=[l]{\#}
}
\lstdefinelanguage{python}{
  morekeywords={from,import,def,class,return,if,elif,else,for,in,while,try,except,raise,with,as,not,and,or,is,lambda,yield,True,False,None,await,async},
  sensitive=true,
  morecomment=[l]{\#},
  morestring=[b]',
  morestring=[b]"
}
\lstdefinelanguage{sql}{
  morekeywords={SELECT,FROM,WHERE,INSERT,INTO,VALUES,UPDATE,SET,DELETE,CREATE,TABLE,PRIMARY,KEY,FOREIGN,REFERENCES,UNIQUE,NOT,NULL,INDEX,ON,ALTER,DROP,CONSTRAINT},
  sensitive=false,
  morecomment=[l]{--}
}
\lstdefinelanguage{json}{
  morecomment=[l]{//},
  morestring=[b]",
  showstringspaces=false
}
\lstdefinelanguage{typescript}{
  morekeywords={import,export,from,const,let,var,function,return,if,else,for,of,in,while,switch,case,break,async,await,try,catch,finally,throw,new,class,this,extends,implements,interface,type,enum,readonly,public,private,protected,static,get,set,true,false,null,undefined,void,string,number,boolean,any,unknown,never,typeof,keyof,as},
  sensitive=true,
  morecomment=[l]{//},
  morecomment=[s]{/*}{*/},
  morestring=[b]',
  morestring=[b]",
  morestring=[b]`
}
\lstalias{javascript}{typescript}
\lstset{
  backgroundcolor=\color{codebg},
  basicstyle=\ttfamily\footnotesize,
  frame=single,
  rulecolor=\color{codeframeclr},
  frameround=tttt,
  breaklines=true,
  tabsize=2,
  showstringspaces=false,
  upquote=true,
  columns=fullflexible
}

% ---------- Helper macros ----------
\newcommand{\codefile}[1]{\texttt{#1}}
\newcommand{\kbd}[1]{\texttt{#1}}

% ---------- Code frame environment ----------
% Automatski dodaje fragile i shrink=5, koristi se umjesto \begin{frame}
% Primjer: \begin{codeframe}{Naslov} ... \end{codeframe}
% Primjer s opcijama: \begin{codeframe}[shrink=15]{Naslov} ... \end{codeframe}
\newenvironment{codeframe}[2][]
  {\begin{frame}[fragile,shrink=5,#1]{#2}}
  {\end{frame}}

% ============================================================

\documentclass[aspectratio=149,11pt]{beamer}
\usetheme{CambridgeUS}

\usepackage{fontspec}              % radi samo sa XeLaTeX/LuaLaTeX
\setsansfont{Latin Modern Sans}    % sans-serif font
\setmonofont{Latin Modern Mono}    % monospace font za kod
\usepackage[croatian]{babel}
\usepackage{csquotes}              % \enquote{} → hrvatski navodnici „..."
\usepackage{lmodern}
\usepackage{booktabs}
\usepackage{listings}
\usepackage{xcolor}
\usepackage{hyperref}

% ---------- Autor i logo ----------
\author{mag.ing. Josip Torić (SIT UNIZD)}
\titlegraphic{\includegraphics[height=2.5cm]{../common/SIT_logo_tekst1.pdf}}

% ---------- Listings ----------
\definecolor{codebg}{RGB}{248,248,248}
\definecolor{codeframeclr}{RGB}{220,220,220}
\lstdefinelanguage{bash}{
  morekeywords={sudo,apt,brew,pip,pipx,uv,winget,python,git,flask,source,deactivate,code,gunicorn,docker,compose,heroku,render,railway,node,npm,pnpm,yarn,uvicorn,alembic,pytest,ruff,black,precommit},
  sensitive=false,
  morecomment=[l]{\#}
}
\lstdefinelanguage{python}{
  morekeywords={from,import,def,class,return,if,elif,else,for,in,while,try,except,raise,with,as,not,and,or,is,lambda,yield,True,False,None,await,async},
  sensitive=true,
  morecomment=[l]{\#},
  morestring=[b]',
  morestring=[b]"
}
\lstdefinelanguage{sql}{
  morekeywords={SELECT,FROM,WHERE,INSERT,INTO,VALUES,UPDATE,SET,DELETE,CREATE,TABLE,PRIMARY,KEY,FOREIGN,REFERENCES,UNIQUE,NOT,NULL,INDEX,ON,ALTER,DROP,CONSTRAINT},
  sensitive=false,
  morecomment=[l]{--}
}
\lstdefinelanguage{json}{
  morecomment=[l]{//},
  morestring=[b]",
  showstringspaces=false
}
\lstdefinelanguage{typescript}{
  morekeywords={import,export,from,const,let,var,function,return,if,else,for,of,in,while,switch,case,break,async,await,try,catch,finally,throw,new,class,this,extends,implements,interface,type,enum,readonly,public,private,protected,static,get,set,true,false,null,undefined,void,string,number,boolean,any,unknown,never,typeof,keyof,as},
  sensitive=true,
  morecomment=[l]{//},
  morecomment=[s]{/*}{*/},
  morestring=[b]',
  morestring=[b]",
  morestring=[b]`
}
\lstalias{javascript}{typescript}
\lstset{
  backgroundcolor=\color{codebg},
  basicstyle=\ttfamily\footnotesize,
  frame=single,
  rulecolor=\color{codeframeclr},
  frameround=tttt,
  breaklines=true,
  tabsize=2,
  showstringspaces=false,
  upquote=true,
  columns=fullflexible
}

% ---------- Helper macros ----------
\newcommand{\codefile}[1]{\texttt{#1}}
\newcommand{\kbd}[1]{\texttt{#1}}

% ---------- Code frame environment ----------
% Automatski dodaje fragile i shrink=5, koristi se umjesto \begin{frame}
% Primjer: \begin{codeframe}{Naslov} ... \end{codeframe}
% Primjer s opcijama: \begin{codeframe}[shrink=15]{Naslov} ... \end{codeframe}
\newenvironment{codeframe}[2][]
  {\begin{frame}[fragile,shrink=5,#1]{#2}}
  {\end{frame}}


\title[Razvoj web aplikacija]{Razvoj web aplikacija - Predavanje 7 \\ Frontend: arhitektura SPA, router i layout}
\date{23.11.2026.}

\begin{document}

\frame{\titlepage}

% ------------------------------------------------------------
\begin{frame}{Plan predavanja}
\begin{enumerate}
  \item Što je SPA i zašto ne radimo MPA
  \item Vue 3 i Composition API
  \item Vite: dev server i build pipeline
  \item Postavljanje projekta (scaffold)
  \item Vue Router: rute, parametri, nested rute
  \item Layout i navigacija
  \item Stilizacija: scoped styles i globalni CSS
  \item TypeScript: setup i integracija s Vue 3
  \item Commitovi za ovo predavanje
\end{enumerate}

\vspace{0.6em}
Na kraju predavanja: Vue projekt s routingom, layoutom i stubiranim stranicama za admin i club korisnika.
\end{frame}

% ------------------------------------------------------------
\begin{frame}{Ishodi učenja}
Nakon ovog predavanja student može:
\begin{itemize}
  \item objasniti razliku između SPA i klasične server-rendered aplikacije
  \item kreirati Vue 3 projekt s Vite-om i TypeScriptom
  \item definirati rute s parametrima i nested strukturom u Vue Routeru
  \item implementirati zajednički layout s navigacijom
  \item organizirati frontend projekt po mapama (views, components, router, stores, api)
  \item tipizirati props, Pinia store i Axios pozive TypeScriptom
\end{itemize}
\end{frame}

% ------------------------------------------------------------
\begin{frame}{SPA vs MPA}
\begin{columns}[T]
\column{0.48\textwidth}
\textbf{MPA (Multi-Page Application)}
\begin{itemize}
  \item Server generira HTML za svaku stranicu
  \item Svaki klik = novi HTTP request + full page reload
  \item Backend i frontend su spojeni (Django templates, Jinja2, PHP)
  \item Jednostavniji SEO
\end{itemize}

\column{0.48\textwidth}
\textbf{SPA (Single-Page Application)}
\begin{itemize}
  \item Browser učita jedan HTML + JS bundle
  \item Navigacija mijenja DOM bez reloada
  \item Frontend i backend komuniciraju kroz API (JSON)
  \item Brži UX nakon inicijalnog loada
\end{itemize}
\end{columns}

\vspace{0.6em}
Mi radimo SPA jer već imamo REST API. Frontend je potpuno odvojen od backenda - može ga razvijati drugi tim.
\end{frame}

% ------------------------------------------------------------
\begin{frame}{Kako SPA zapravo radi}
\begin{enumerate}
  \item Browser dohvati \texttt{index.html} - jedini HTML koji server ikad pošalje
  \item HTML učita JavaScript bundle - entry point je \texttt{src/main.ts}
  \item JS preuzima renderiranje: kreira DOM, sluša eventove, mijenja prikaz
  \item Klik na link ne šalje request serveru - JS mijenja URL i renderira novu komponentu
  \item Podaci dolaze asinkrono kroz API pozive (JSON)
\end{enumerate}

\vspace{0.4em}
Server nakon inicijalnog loada služi samo kao API. Cijeli UI živi u browseru.
\end{frame}

% ------------------------------------------------------------
\begin{frame}{Prednosti odvojenog frontenda}
\begin{tabular}{@{}lp{10cm}@{}}
  \toprule
  \textbf{Prednost} & \textbf{Objašnjenje} \\
  \midrule
  Nezavisni deploy & Frontend i backend se deployaju zasebno. Promjena boje gumba ne zahtijeva deploy backenda. \\
  Različiti timovi & Frontend i backend tim rade paralelno. Dogovor je API ugovor (Swagger). \\
  API-first & Isti backend služi web, mobilnu aplikaciju i integracije. \\
  Tehnološka sloboda & Frontend može biti Vue, React ili mobilna app - backend ne zna i ne mari. \\
  \bottomrule
\end{tabular}
\end{frame}

% ------------------------------------------------------------
\begin{frame}{Reaktivni sustav: koncept}
Vue koristi \textbf{reaktivni sustav} koji automatski prati zavisnosti između podataka i UI-ja.

\vspace{0.6em}
\begin{center}
\texttt{podatak se promijeni $\rightarrow$ Vue detektira $\rightarrow$ DOM se ažurira}
\end{center}

\vspace{0.6em}
\begin{itemize}
  \item Nema ručnog \texttt{document.getElementById().innerHTML = ...}
  \item Vue interno koristi JavaScript \texttt{Proxy} za praćenje pristupa i promjena
  \item Kad čitaš varijablu u templateu, Vue zabilježi tu zavisnost
  \item Kad se varijabla promijeni, Vue zna točno koje dijelove DOM-a treba ažurirati
\end{itemize}

\vspace{0.4em}
Ovo je ključna razlika od jQuery/vanilla JS pristupa - ne razmišljate o DOM-u, razmišljate o podacima.
\end{frame}

% ------------------------------------------------------------
\begin{frame}{Virtual DOM}
Vue ne mijenja stvarni DOM direktno. Koristi \textbf{Virtual DOM} - laganu JavaScript reprezentaciju UI-ja.

\vspace{0.6em}
\begin{enumerate}
  \item Komponenta renderira Virtual DOM stablo (JS objekti)
  \item Kad se podatak promijeni, renderira se novo Virtual DOM stablo
  \item Vue uspoređuje staro i novo stablo (diffing)
  \item Samo razlike se primijene na stvarni DOM
\end{enumerate}

\vspace{0.6em}
Stvarni DOM je spor - svaka promjena može triggerirati layout recalculation, repaint i reflow. Virtual DOM minimizira broj stvarnih DOM operacija.
\end{frame}

% ------------------------------------------------------------
\begin{frame}{Vue 3: zašto baš Vue}
\begin{itemize}
  \item Manji learning curve od Reacta - template syntax je bliži HTML-u
  \item Composition API daje fleksibilnost sličnu React Hooksima
  \item Ekosustav: Vue Router, Pinia, Vite - sve službeno i dobro integrirano
  \item Odlična dokumentacija na \texttt{vuejs.org}
\end{itemize}

\vspace{0.6em}
Vue koristi \textbf{reaktivni sustav} - kad se podatak promijeni, UI se automatski ažurira. Nema ručnog DOM manipuliranja.
\end{frame}

% ------------------------------------------------------------
\begin{frame}{Composition API vs Options API}
Vue 3 nudi dva načina pisanja komponenti. Mi koristimo \textbf{Composition API} jer je fleksibilniji i bolje skalira.

\vspace{0.4em}
\begin{columns}[T]
\column{0.48\textwidth}
\textbf{Options API}
\begin{itemize}
  \item \texttt{data()}, \texttt{methods}, \texttt{computed}, \texttt{watch}
  \item Logika raspoređena po sekcijama
  \item Teško za ponovnu upotrebu logike
\end{itemize}

\column{0.48\textwidth}
\textbf{Composition API}
\begin{itemize}
  \item \texttt{ref()}, \texttt{computed()}, \texttt{watch()}
  \item Logika grupirana po značenju
  \item Composables za dijeljenje logike
  \item \texttt{<script setup>} sintaksa
\end{itemize}
\end{columns}
\end{frame}

% ------------------------------------------------------------
\begin{frame}{ref() vs reactive()}
\begin{tabular}{@{}lp{5cm}p{5cm}@{}}
  \toprule
  & \textbf{ref()} & \textbf{reactive()} \\
  \midrule
  Tip & Primitivi i objekti & Samo objekti \\
  Pristup & \texttt{count.value} u scriptu, \texttt{count} u templateu & Direktno: \texttt{state.count} \\
  Reassign & Da (\texttt{count.value = 5}) & Ne (gubi reaktivnost) \\
  Destrukturiranje & Sigurno s \texttt{toRefs()} & Gubi reaktivnost \\
  \bottomrule
\end{tabular}

\vspace{0.6em}
Preporuka: koristite \texttt{ref()} za sve. Jednostavniji je, radi s primitivima i objektima, i \texttt{.value} čini eksplicitnim da radite s reaktivnom varijablom.
\end{frame}

% ------------------------------------------------------------
\begin{frame}{Životni ciklus komponente}
Svaka Vue komponenta prolazi kroz definirane faze:

\vspace{0.4em}
\begin{center}
\texttt{kreiranje $\rightarrow$ montiranje $\rightarrow$ ažuriranje $\rightarrow$ uništavanje}
\end{center}

\vspace{0.6em}
\begin{tabular}{@{}ll@{}}
  \toprule
  \textbf{Hook} & \textbf{Kad se poziva} \\
  \midrule
  \texttt{onMounted} & Komponenta je renderirana u DOM. API pozivi idu ovdje. \\
  \texttt{onUpdated} & Reaktivni podatak se promijenio, DOM je ažuriran. \\
  \texttt{onUnmounted} & Komponenta je uklonjena. Cleanup (intervali, listeneri). \\
  \texttt{onBeforeUnmount} & Neposredno prije uklanjanja. \\
  \bottomrule
\end{tabular}

\vspace{0.4em}
\texttt{onMounted} ćete koristiti najčešće - za dohvat podataka kad se stranica učita.
\end{frame}

% ------------------------------------------------------------
\begin{frame}{Template sintaksa: direktive}
Vue direktive su specijalni atributi koji počinju s \texttt{v-}:

\vspace{0.4em}
\begin{tabular}{@{}lp{8cm}@{}}
  \toprule
  \textbf{Direktiva} & \textbf{Značenje} \\
  \midrule
  \texttt{v-if} / \texttt{v-else} & Uvjetno renderiranje - element ne postoji u DOM-u \\
  \texttt{v-show} & Uvjetno sakrivanje - element postoji ali je \texttt{display: none} \\
  \texttt{v-for} & Iteracija kroz array ili objekt \\
  \texttt{v-bind} (\texttt{:}) & Dinamički atribut - \texttt{:class}, \texttt{:href}, \texttt{:disabled} \\
  \texttt{v-on} (\texttt{@}) & Event listener - \texttt{@click}, \texttt{@submit}, \texttt{@input} \\
  \texttt{v-model} & Dvosmjerni binding - spaja input s varijablom \\
  \bottomrule
\end{tabular}

\vspace{0.4em}
\texttt{v-if} uklanja element iz DOM-a (skuplji toggle, jeftiniji init). \texttt{v-show} ga samo skriva (jeftiniji toggle, skuplji init).
\end{frame}

% ------------------------------------------------------------
\begin{frame}{Props i Emits}
Parent-child komunikacija u Vueu je \textbf{jednosmjerna}:

\vspace{0.6em}
\begin{center}
\texttt{Parent $\xrightarrow{\text{props}}$ Child $\xrightarrow{\text{emit}}$ Parent}
\end{center}

\vspace{0.6em}
\begin{itemize}
  \item \textbf{Props} - podaci teku od parenta prema childu. Child ih nikad ne smije mijenjati.
  \item \textbf{Emits} - child šalje event parentu kad se nešto dogodi (klik, submit, promjena).
  \item Parent odlučuje što napraviti s eventom - child ne zna.
\end{itemize}

\vspace{0.6em}
\texttt{defineProps()} deklarira koje propse komponenta prima. \texttt{defineEmits()} deklarira koje evente emitira. Oboje su compile-time makroi u \texttt{<script setup>}.
\end{frame}

% ------------------------------------------------------------
\begin{frame}{Computed vs Watch}
\begin{columns}[T]
\column{0.48\textwidth}
\textbf{computed()}
\begin{itemize}
  \item Derivirani podatak iz postojećih
  \item Kešira se - recalculira samo kad se zavisnosti promijene
  \item Read-only (uglavnom)
  \item Primjer: filtrirani popis, ukupna cijena
\end{itemize}

\column{0.48\textwidth}
\textbf{watch()}
\begin{itemize}
  \item Reagira na promjenu podatka
  \item Pokreće side-effect (API poziv, log)
  \item Ne kešira ništa
  \item Primjer: dohvat podataka kad se ruta promijeni
\end{itemize}
\end{columns}

\vspace{0.6em}
Ako možeš izraziti nešto kao \texttt{computed}, koristi \texttt{computed}. \texttt{watch} je za situacije kad promjena podatka treba pokrenuti akciju van renderiranja.
\end{frame}

% ------------------------------------------------------------
\begin{frame}{Slots: kompozicija komponenti}
Slots omogućuju parentu da \enquote{ubaci} sadržaj unutar child komponente:

\vspace{0.4em}
\begin{tabular}{@{}lp{8cm}@{}}
  \toprule
  \textbf{Tip slota} & \textbf{Upotreba} \\
  \midrule
  Default slot & Jedan blok sadržaja - \texttt{<MyCard>tekst</MyCard>} \\
  Named slots & Više blokova - \texttt{<template \#header>}, \texttt{<template \#footer>} \\
  Scoped slots & Child pruža podatke parentu za renderiranje \\
  \bottomrule
\end{tabular}

\vspace{0.6em}
Slots su Vue verzija \enquote{composition over inheritance}. Umjesto da child komponenta zna sve o sadržaju koji prikazuje, parent šalje sadržaj, a child ga pozicionira.
\end{frame}

% ------------------------------------------------------------
\begin{frame}[fragile,shrink=5]{Prva Vue komponenta}
\begin{lstlisting}[language=html]
<script setup lang="ts">
import { ref, computed } from 'vue'

const count = ref(0)
const doubled = computed(() => count.value * 2)

function increment() {
  count.value++
}
</script>

<template>
  <button @click="increment">
    Klikova: {{ count }} (duplo: {{ doubled }})
  </button>
</template>
\end{lstlisting}

\vspace{0.4em}
\texttt{ref()} stvara reaktivnu varijablu. U scriptu se pristupa s \texttt{.value}, u templateu direktno.
\end{frame}

% ------------------------------------------------------------
\begin{frame}[fragile,shrink=5]{Scaffold projekta s TypeScriptom}
\begin{lstlisting}[language=bash]
# Interaktivni wizard - odaberi TypeScript: Yes
npm create vue@latest web

cd web
npm install

# Dodaj preostale dependencije
npm install axios

# Type-check bez bundlanja
npm run type-check   # vue-tsc --noEmit

# Dev server s hot-reload
npm run dev    # http://localhost:5173
\end{lstlisting}

\vspace{0.4em}
Wizard konfigurira \texttt{tsconfig.json}, \texttt{env.d.ts} i Vue Router + Pinia po izboru. Sve \texttt{.vue} datoteke automatski dobivaju \texttt{<script setup lang="ts">}.
\end{frame}

% ------------------------------------------------------------
\begin{frame}{TypeScript u Vue projektu: zašto}
\begin{itemize}
  \item \textbf{Greške pri pisanju, ne pri izvršavanju} - TypeScript javlja pogrešan tip prije nego što browser uopće pokrene kod
  \item \textbf{Autocomplete koji zna kontekst} - IDE prikazuje dostupna polja na API responseu, propse komponente, metode storea
  \item \textbf{Sigurni refactoring} - preimenuješ interface, TypeScript pokaže svako mjesto koje treba popraviti
  \item \textbf{Vue 3 je pisan u TypeScriptu} - podrška je ugrađena od početka, nema workaroundova ni zastarjelih tipova
\end{itemize}

\vspace{0.6em}
TypeScript ne usporava razvoj na ovom projektu - ubrzava ga jer eliminira cijelu kategoriju runtime bugova (undefined is not a function, pogrešan naziv polja).
\end{frame}

% ------------------------------------------------------------
\begin{frame}[fragile,shrink=5]{Tipiziranje komponenti: defineProps i defineEmits}
\begin{lstlisting}[language=typescript]
// src/components/FormField.vue
<script setup lang="ts">

// Genericki interface - bez runtime overhead
defineProps<{
  label: string
  error?: string
  modelValue: string | number
}>()

defineEmits<{
  'update:modelValue': [value: string]
}>()

// Domain objekt - dijeli se s Pinia storeom
interface Lifter {
  id: number
  first_name: string
  last_name: string
  gender: 'M' | 'F'
}
</script>
\end{lstlisting}
\end{frame}

% ------------------------------------------------------------
\begin{frame}[fragile,shrink=5]{Tipiziranje Pinia storea}
\begin{lstlisting}[language=typescript]
// src/stores/lifters.ts
import { defineStore } from 'pinia'
import { ref } from 'vue'
import type { Lifter } from '@/types'

export const useLifterStore = defineStore('lifters', () => {
  const lifters = ref<Lifter[]>([])
  const loading = ref(false)
  const error = ref<string | null>(null)

  async function fetchLifters(clubId: number): Promise<void> {
    loading.value = true
    // ...
  }

  return { lifters, loading, error, fetchLifters }
})
\end{lstlisting}

\vspace{0.4em}
\texttt{useLifterStore()} vraća potpuno tipiziran objekt - autocomplete radi za cijeli store.
\end{frame}

% ------------------------------------------------------------
\begin{frame}[fragile,shrink=5]{Tipiziranje Axios poziva}
\begin{lstlisting}[language=typescript]
// src/api/lifters.ts
import api from './client'
import type { Lifter } from '@/types'

export async function getLifters(clubId: number): Promise<Lifter[]> {
  const { data } = await api.get<Lifter[]>(
    `/clubs/${clubId}/lifters`
  )
  return data   // TypeScript zna: data je Lifter[]
}

export async function createLifter(
  clubId: number,
  body: Omit<Lifter, 'id'>
): Promise<Lifter> {
  const { data } = await api.post<Lifter>(
    `/clubs/${clubId}/lifters`, body
  )
  return data
}
\end{lstlisting}
\end{frame}

% ------------------------------------------------------------
\begin{frame}{Centralni types/ folder}
Interfacei se definiraju jednom, koriste svugdje:

\vspace{0.4em}
\begin{tabular}{@{}lp{9cm}@{}}
  \toprule
  \textbf{File} & \textbf{Sadržaj} \\
  \midrule
  \texttt{src/types/index.ts} & \texttt{Lifter}, \texttt{Club}, \texttt{Competition}, \texttt{Registration} \\
  \texttt{src/types/auth.ts} & \texttt{User}, \texttt{TokenPair}, \texttt{LoginPayload} \\
  \texttt{src/types/api.ts} & \texttt{PaginatedResponse<T>}, \texttt{ApiError} \\
  \bottomrule
\end{tabular}

\vspace{0.6em}
Preslikavanje s Pydantic schemom je 1:1 - isti nazivi polja, isti tipovi. Kad backend promijeni schemu, TypeScript u frontendu odmah pokaže gdje se koristi staro polje. Ovo je prednost API-first pristupa.
\end{frame}

% ------------------------------------------------------------
\begin{frame}{tsconfig: relevantne opcije}
\begin{tabular}{@{}lp{8cm}@{}}
  \toprule
  \textbf{Opcija} & \textbf{Što radi} \\
  \midrule
  \texttt{strict: true} & Sve provjere uključene (preporučeno za novi projekt) \\
  \texttt{paths: \{"@/*": ["./src/*"]\}} & Alias \texttt{@/} za \texttt{src/} - mora biti isti kao u \texttt{vite.config.ts} \\
  \texttt{include: ["src/**/*"]} & TypeScript gleda samo naš kod, ne \texttt{node\_modules} \\
  \texttt{vue-tsc} & Type-check \texttt{.vue} datoteka - \texttt{tsc} sam ne razumije Vue SFC \\
  \bottomrule
\end{tabular}

\vspace{0.6em}
\texttt{npm run type-check} pokreće \texttt{vue-tsc --noEmit} - provjerava tipove bez generiranja outputa. Dodajte ga u CI pipeline uz lint i testove.
\end{frame}

% ------------------------------------------------------------
\begin{frame}{Vite vs Webpack}
\begin{tabular}{@{}lp{5cm}p{5cm}@{}}
  \toprule
  & \textbf{Webpack} & \textbf{Vite} \\
  \midrule
  Dev start & Bundla cijeli projekt prije starta (sporo) & Servira module nativno (instant) \\
  HMR & Cijeli modul se zamijeni & Samo promijenjeni modul \\
  Konfiguracija & Kompleksna (\texttt{webpack.config.js}) & Minimalna (\texttt{vite.config.ts}) \\
  Prod build & Webpack bundler & Rollup bundler \\
  \bottomrule
\end{tabular}

\vspace{0.6em}
Vite koristi činjenicu da moderni browseri podržavaju ES modules nativno. U developmentu nema bundlanja - svaki file se servira pojedinačno. Rezultat: dev server se pokrene za sekunde umjesto minuta.
\end{frame}

% ------------------------------------------------------------
\begin{frame}{Struktura frontend projekta}
\begin{columns}[T]
\column{0.48\textwidth}
\codefile{web/src/}
\begin{itemize}
  \item \codefile{views/} - stranice (jedna po ruti)
  \item \codefile{components/} - reusable UI dijelovi
  \item \codefile{router/} - definicija ruta
  \item \codefile{stores/} - Pinia storeovi
  \item \codefile{api/} - HTTP client i pozivi
  \item \codefile{App.vue} - root komponenta
  \item \codefile{main.ts} - entry point
\end{itemize}

\column{0.48\textwidth}
\textbf{Konvencije:}
\begin{itemize}
  \item Views su stranice - \texttt{LoginView.vue}, \texttt{DashboardView.vue}
  \item Components su dijelovi - \texttt{NavBar.vue}, \texttt{LoadingSpinner.vue}
  \item Jedan view = jedna ruta
  \item Store po domeni - \texttt{authStore}, \texttt{lifterStore}
\end{itemize}
\end{columns}
\end{frame}

% ------------------------------------------------------------
\begin{frame}[fragile,shrink=5]{Vue Router: osnovna konfiguracija}
\begin{lstlisting}[language=typescript]
// src/router/index.ts
import { createRouter, createWebHistory } from 'vue-router'
import LoginView from '@/views/LoginView.vue'
import DashboardView from '@/views/DashboardView.vue'

const routes = [
  { path: '/login', name: 'login', component: LoginView },
  {
    path: '/',
    name: 'dashboard',
    component: DashboardView,
  },
]

const router = createRouter({
  history: createWebHistory(),
  routes,
})

export default router
\end{lstlisting}
\end{frame}

% ------------------------------------------------------------
\begin{frame}[fragile,shrink=5]{Rute s parametrima i nested rute}
\begin{lstlisting}[language=typescript]
// Ruta s parametrom
{ path: '/clubs/:clubId/lifters', name: 'lifters',
  component: LiftersView }

// Pristup parametru u komponenti
const route = useRoute()
const clubId = route.params.clubId

// Nested rute - dijele layout
{
  path: '/admin',
  component: AdminLayout,
  children: [
    { path: 'clubs', component: AdminClubsView },
    { path: 'competitions', component: AdminCompetitionsView },
  ],
}
\end{lstlisting}

\vspace{0.4em}
Nested rute koriste \texttt{<RouterView />} unutar parent komponente za renderiranje child rute.
\end{frame}

% ------------------------------------------------------------
\begin{frame}{Hash mode vs History mode}
Vue Router podržava dva načina upravljanja URL-ovima:

\vspace{0.4em}
\begin{tabular}{@{}lp{5cm}p{5cm}@{}}
  \toprule
  & \textbf{Hash mode} & \textbf{History mode} \\
  \midrule
  URL & \texttt{app.com/\#/clubs} & \texttt{app.com/clubs} \\
  Kako radi & Browser ignorira \# dio & HTML5 History API \\
  Server config & Ne treba & Treba fallback na index.html \\
  SEO & Lošije (\# se ne indexira) & Bolje \\
  \bottomrule
\end{tabular}

\vspace{0.6em}
Mi koristimo \texttt{createWebHistory()} - history mode. URL-ovi izgledaju prirodno, ali server mora sve nepoznate rute preusmjeriti na \texttt{index.html} jer je routing na klijentu.
\end{frame}

% ------------------------------------------------------------
\begin{frame}{Lazy loading ruta}
Kad aplikacija naraste, ne želimo učitati sav JavaScript odjednom:

\vspace{0.4em}
\begin{itemize}
  \item \textbf{Eager loading}: sve komponente se učitaju pri startu - sporiji inicijalni load
  \item \textbf{Lazy loading}: komponenta se učita tek kad korisnik navigira na tu rutu
\end{itemize}

\vspace{0.6em}
Umjesto \texttt{import AdminView from '@/views/AdminView.vue'} koristimo:

\vspace{0.2em}
\texttt{component: () => import('@/views/AdminView.vue')}

\vspace{0.6em}
Vite automatski radi \textbf{code splitting} - svaki lazy-loaded view postaje zasebni JS chunk. Korisnik preuzima samo kod za stranicu koju gleda.
\end{frame}

% ------------------------------------------------------------
\begin{frame}[fragile,shrink=5]{Layout s navigacijom}
\begin{lstlisting}[language=html]
<!-- src/App.vue -->
<script setup lang="ts">
import NavBar from '@/components/NavBar.vue'
</script>

<template>
  <NavBar />
  <main class="container">
    <RouterView />
  </main>
</template>
\end{lstlisting}

\vspace{0.4em}
\texttt{<RouterView />} je placeholder - Vue Router renderira komponentu koja odgovara trenutnoj ruti. Navigacija (\texttt{NavBar}) ostaje fiksna jer je izvan \texttt{<RouterView />}.
\end{frame}

% ------------------------------------------------------------
\begin{frame}[fragile,shrink=5]{Stilizacija: scoped styles u SFC-u}
Treći blok svake .vue datoteke je \texttt{<style>}. S atributom \texttt{scoped} stilovi vrijede samo za tu komponentu:

\begin{lstlisting}[language=html]
<!-- src/components/NavBar.vue -->
<style scoped>
nav {
  display: flex;
  gap: 1rem;
  padding: 0.75rem 1rem;
  border-bottom: 1px solid #ddd;
}
.error { color: #b00020; font-size: 0.875rem; }
button:disabled { opacity: 0.5; cursor: not-allowed; }
</style>
\end{lstlisting}

\vspace{0.4em}
Vue svakoj komponenti doda jedinstveni data-atribut i prepiše selektore - zato \texttt{.error} iz NavBara ne curi u formu. Svaka klasa koju koristite u templateu (\texttt{error}, \texttt{badge-green}, \texttt{btn-danger}) mora biti definirana ovdje ili globalno.
\end{frame}

% ------------------------------------------------------------
\begin{frame}[fragile,shrink=5]{Globalni stilovi i responsivnost}
\begin{lstlisting}[language=html]
/* src/assets/main.css - importan u main.ts */
:root { --primary: #1f6feb; }
.container { max-width: 960px;
             margin: 0 auto; padding: 1rem; }
.badge-green { background: #d1f7d6;
               padding: 2px 8px; border-radius: 8px; }

@media (max-width: 600px) {
  .container { padding: 0.5rem; }
  table { display: block; overflow-x: auto; }
}
\end{lstlisting}

\vspace{0.4em}
Globalno ide ono što dijele sve komponente: reset, layout klase, badgevi. Scaffold već stvara \codefile{assets/main.css} - proširujete njega. \texttt{@media} breakpoint čini tablice upotrebljivima na mobitelu.

\vspace{0.4em}
Za projekt je dovoljan vlastiti mali CSS. Komponentne biblioteke (PrimeVue, Vuetify) smijete koristiti, ali ih ne obrađujemo - na obrani morate objasniti i njihov kod.
\end{frame}

% ------------------------------------------------------------
\begin{frame}[fragile,shrink=5]{Navigacija: RouterLink}
\begin{lstlisting}[language=html]
<!-- src/components/NavBar.vue -->
<template>
  <nav>
    <RouterLink to="/">Dashboard</RouterLink>
    <RouterLink to="/clubs">Klubovi</RouterLink>
    <RouterLink :to="`/clubs/${clubId}/lifters`">
      Natjecatelji
    </RouterLink>
    <button @click="logout">Odjava</button>
  </nav>
</template>
\end{lstlisting}

\vspace{0.4em}
\texttt{RouterLink} renderira \texttt{<a>} tag ali presreće klik - umjesto full page reloada, mijenja samo komponentu u \texttt{<RouterView />}.
\end{frame}

% ------------------------------------------------------------
\begin{frame}[fragile,shrink=5]{Programska navigacija}
\begin{lstlisting}[language=typescript]
import { useRouter } from 'vue-router'

const router = useRouter()

// Navigacija po imenu rute
router.push({ name: 'lifters', params: { clubId: 5 } })

// Navigacija po pathu
router.push('/login')

// Zamjena (ne dodaje u history)
router.replace('/dashboard')

// Natrag
router.back()
\end{lstlisting}

\vspace{0.4em}
\texttt{router.push} dodaje u browser history (back button radi). \texttt{router.replace} zamjenjuje trenutni entry - koristi se nakon logina da korisnik ne može ići \enquote{back} na login formu.
\end{frame}

% ------------------------------------------------------------
\begin{frame}{Konvencije i alati}
\textbf{Konvencije imenovanja:}

\vspace{0.2em}
\begin{tabular}{@{}llp{5.5cm}@{}}
  \toprule
  \textbf{Što} & \textbf{Konvencija} & \textbf{Primjer} \\
  \midrule
  Komponenta & PascalCase & \texttt{NavBar.vue}, \texttt{LifterForm.vue} \\
  Props / emiti & camelCase / kebab-case & \texttt{firstName}, \texttt{@item-click} \\
  Composables & \texttt{use} prefix & \texttt{useAuth()}, \texttt{useLifters()} \\
  Store & \texttt{use...Store} & \texttt{useAuthStore()} \\
  \bottomrule
\end{tabular}

\vspace{0.8em}
\textbf{Vue DevTools} - browser ekstenzija, obavezni alat za razvoj:
\begin{itemize}
  \item \textbf{Components} - stablo komponenti, props i lokalni state
  \item \textbf{Pinia} - svi storeovi, njihov state, mogućnost uređivanja
  \item \textbf{Routes} - trenutna ruta, parametri, matched records
\end{itemize}

\vspace{0.4em}
Kad nešto ne radi - otvorite Vue DevTools. Vidjet ćete je li komponenta renderirana, imaju li propsi ispravne vrijednosti i je li store ažuriran.
\end{frame}

% ------------------------------------------------------------
\begin{frame}{Česte greške početnika u Vueu}
\begin{enumerate}
  \item \textbf{Mutiranje propsa} - child mijenja prop direktno. Vue javlja warning, ali podatak se ne ažurira u parentu. Koristiti emit.
  \item \textbf{Zaboravljen :key u v-for} - bez keya Vue ne može efikasno pratiti elemente. Koristiti \texttt{:key="item.id"}, nikad index.
  \item \textbf{Reaktivnost na nested objektima} - \texttt{ref({a: {b: 1}})} - promjena \texttt{obj.value.a.b} se detektira jer Vue 3 koristi Proxy, ali zamjena cijelog objekta zahtijeva \texttt{.value = \{...\}}.
  \item \textbf{Zaboravljen .value} - u \texttt{<script setup>} pristupate \texttt{ref} s \texttt{.value}, u templateu bez njega. Miješanje dovodi do bugova.
\end{enumerate}
\end{frame}

% ------------------------------------------------------------
\begin{frame}[fragile,shrink=5]{Registracija routera i Pinie u main.ts}
\begin{lstlisting}[language=typescript]
// src/main.ts
import { createApp } from 'vue'
import { createPinia } from 'pinia'
import App from './App.vue'
import router from './router'

const app = createApp(App)
app.use(createPinia())
app.use(router)
app.mount('#app')
\end{lstlisting}

\vspace{0.4em}
\texttt{app.use()} registrira plugine. Redoslijed je bitan - Pinia prije routera jer route guardovi mogu koristiti storeove.
\end{frame}

% ------------------------------------------------------------
\begin{frame}{Commitovi za Predavanje 7}
\begin{itemize}
  \item \texttt{feat(web): scaffold app layout, router, and navigation shell}
  \item \texttt{feat(web): add route structure for admin and club areas}
  \item \texttt{chore(web): add eslint/prettier and basic project conventions}
  \item \texttt{docs: add frontend structure and component guidelines}
\end{itemize}

\vspace{0.4em}
\end{frame}

% ------------------------------------------------------------
\begin{frame}{Što ide sljedeće (Predavanje 8)}
\begin{itemize}
  \item API komunikacija s backendom (Axios/fetch)
  \item Auth flow: login forma, token storage, refresh
  \item Route guardovi za zaštitu ruta po roli
  \item Error handling na frontendu
\end{itemize}

\vspace{0.8em}
\begin{center}
\Large Pitanja?
\end{center}
\end{frame}

\end{document}
